% Рассматривается задача оценки перемещения мобильного робота, конструкция которого
% предполагает наличие одновременно стерео-камеры и лидара. Требуется исследовать и совместить
% методы визуальной и лидарной одометрии для повышения устойчивости определения
% траектории движения по сравнению с использованием только одного из датчиков. Базовых
% метод метод совмещения данных должен быть выбран по результатам предварительного анализа. 
% Разработанный прототип необходимо протестировать на открытых наборах данных, 
% в частности на The KITTI Vision Benchmark Suite (Visual Odometry).
% Цель работы — разработать и исследовать мультимодальный алгоритм одометрии,
% комбинирующий информацию от стерео-камеры и лидара, и определить рекомендации по
% применению такого подхода в робототехнике.
% Задачи:
% 1. Проведение аналитического обзора в области методов одометрии на основе данных
% камеры и лидара, включая способы их комбинирования.
% 2. Выбор подходящего метода интеграции визуальных и лидарных данных.
% 3. Реализация прототипа мультимодального алгоритма одометрии, обеспечивающего
% оценку движения мобильного робота.
% 4. Проведение серии экспериментов на открытых датасетах, сравнительного анализа
% полученных результатов с раздельными схемами (только камера / только лидар), а также
% другими представителями мультимодальных методов.
% 5. Формирование методических рекомендаций по использованию и настройке полученного
% решения, включая потенциальные сценарии и ограничения.
% Требования к разрабатываемому решению:
% 1. Разрабатываемая система рассчитана на использование данных стерео RGB или монохромной
% камеры с разрешением 1382 x 512 и частотой не менее 10 Гц, а также 3D лидара с 64
% вертикальными каналами, работающем на частоте не менее 10 Гц.
% 2. Подразумеваются доступными данные калибровки, включающие в себя параметры
% камеры, а также взаимное расположение камеры и лидара.
% 3. Средняя частота обработки данных должна составлять не менее 8 Гц.
% 4. По результатам серии экспериментов средняя относительная ошибка оценки
% пройденного пути не должна превышать 2%, а относительная ошибка ориентации – 0.006
% град/м на The KITTI Vision Benchmark Suite (Visual Odometry).
% 5. Аппаратные ограничения: объём доступной оперативной памяти не менее 8 ГБ, CPU с
% частотой 2 ГГц и выше, 8 ГБ видеопамяти.
% 6. Наличие описания процесса настройки алгоритма, шагов калибровки и воспроизведения
% экспериментов.
% Исходные данные:
% 1. Lee D. et al. LiDAR odometry survey: recent advancements and remaining challenges
% //Intelligent Service Robotics. – 2024. – Т. 17. – No. 2. – С. 95-118.
% 2. He M. et al. A review of monocular visual odometry //The Visual Computer. – 2020. – Т. 36. –
% No. 5. – С. 1053-1065.
% 3. Vizzo I. et al. Kiss-icp: In defense of point-to-point icp–simple, accurate, and robust
% registration if done the right way //IEEE Robotics and Automation Letters. – 2023. – Т. 8. – No. 2.
% – С. 1029-1036.
% 4. Koide K. et al. Glim: 3d range-inertial localization and mapping with gpu-accelerated scan
% matching factors //Robotics and Autonomous Systems. – 2024. – Т. 179. – С. 104750.
% 5. Yuan Z. et al. SDV-LOAM: semi-direct visual–LiDAR Odometry and mapping //IEEE
% Transactions on Pattern Analysis and Machine Intelligence. – 2023. – Т. 45. – No. 9. – С. 11203-
% 11220.
% По результатам разработки должен быть представлен программный прототип
% мультимодального метода одометрии, включающий интегрированную схему работы с
% данными лидара и стерео-камеры, а также отчёт с описанием архитектуры решения,
% сравнительными экспериментами и методическими рекомендациями по использованию.




\newpage
\renewcommand{\contentsname}{\centerline{\large СОДЕРЖАНИЕ}}
\tableofcontents

\newpage
\begin{center}
  \textbf{\large ВВЕДЕНИЕ}
\end{center}
\addcontentsline{toc}{chapter}{ВВЕДЕНИЕ}

В последние годы рынок робототехники и автономных систем стремительно развивается. В частности, растет спрос на мобильных роботов, 
которые всё активнее внедряются в различные отрасли, такие как логистика, сельское хозяйство, городской транспорт и сфера услуг. 
При этом важнейшим вызовом для подобных систем остаётся задача локализации – определения собственного положения и ориентации робота в пространстве. 
Неточности при локализации могут привести к неэффективной или нештатной работе, что в конечном итоге, может стать причиной аварий и сбоев.

Исторически для решения задачи локализации и построения карты (SLAM – Simultaneous Localization and Mapping) используются такие сенсоры, 
как камеры, лидары, IMU (инерционные датчики) и колёсные одометры. Каждая из модальностей имеет ряд преимуществ и недостатков. Например, 
камеры позволяют обрабатывать насыщенную визуальную информацию, что даёт возможность использовать особенности (признаки текстуры, цветные объекты и т. д.) 
и проводить оценку положения на основе сравнительно простого аппаратного обеспечения. Однако в условиях плохого освещения или отсутствия достаточного 
разнообразия признаков эффективность визуальных методов сильно падает. Лидары, напротив, обеспечивают надёжное измерение расстояния 
до объектов в широком диапазоне и с высокой точностью, но не несут в себе цветовой или текстурной информации. Их точность также может падать при определенных погодных условиях (туман, дождь, снегопад).

Именно поэтому возникает необходимость использования мультимодального подхода: в частности, использование одновременно данных камеры и лидара (Visual-LiDAR Odometry) 
обещает более высокую точность и стабильность в самых разных сценариях. Подобные мультимодальные системы уже исследуются и применяются в автомобильной 
отрасли (автономные машины), роботах-курьерах и сервисных платформах, складской логистике и других сферах.

В рамках данной выпускной квалификационной работы ставится цель разработать метод одометрии, совмещающий данные лидара и камеры, а также исследовать его работоспособность. 
Актуальность темы работы обусловлена возросшими требованиями к автономным системам, которым необходимы более точные и 
отказоустойчивые методы оценки собственного движения. Совмещение различных типов сенсорных данных позволяет более полно использовать 
доступную информацию об окружающей среде и снизить риск ошибок в локализации. Научно-практическая значимость работы состоит в том, 
что выводы, сделанные в ходе анализа и экспериментов, могут быть использованы при построении мультимодальных SLAM-систем в различных 
робототехнических приложениях.